% PAGES: 121-126
% Section 3.6 "References" ends before page 121 (owned by a different
% agent's file); the "3.7 Exercises" heading opens directly at the top of
% page 121, so this file begins fresh with that section. All ten exercises
% (pages 121-125) are the whole of Chapter 3's exercise list and are kept
% together in a single \exercises environment so the shared enumerate
% counter is never reset mid-list. Page 126 (PDF page 126, the last page of
% this assignment and of the chapter) is blank -- a transitional verso
% before Chapter 4; see the note at the very end of this file.

\section{Exercises}

\begin{exercises}

\item In three months, the value of an asset with spot price \$50 will be
either \$60 or \$40, with probability one half. What is the value of a
three months at--the--money put on this asset? Assume the risk--free
interest rate is zero.

\item A one period market with four securities and four states at time
$\tau > t_0$ has the following price vector of spot prices at the initial
time $t_0$:
\[
S_{t_0} \;=\; \begin{pmatrix} 4\\ 100\\ 1\\ -2\end{pmatrix},
\]
and the following payoff matrix at time $\tau$:
\[
M_\tau \;=\; \begin{pmatrix} 1 & 1 & 1 & 1\\ 20 & 22 & 24 & 26\\ -5 & -3 & -1 & 2\\ 4 & 2 & 0 & 0\end{pmatrix}.
\]
Find an arbitrage opportunity.

\item A one period market model is made of four securities and has four
states at time $\tau > t_0$. Assume that the market model is complete and
that the state prices are $-1$, $1$, $2$, and $4$, respectively. Find an
arbitrage opportunity.

\item At time $\tau > t_0$, an asset with spot price $S_0$ at time $t_0$
will be worth either $uS_0$, in which case the value of \$1 today would be
$FV_1$, or $dS_0$, in which case the value of \$1 today would be $FV_2$,
with $d < u$. Consider the one period market model with two securities,
i.e., cash and the asset, and two states, i.e., asset value at time $\tau$
equal to $uS_0$ and asset value at time $\tau$ equal to $dS_0$.
\begin{subparts}
\item Show that the payoff matrix at time $\tau$ of this one period market
model is
\[
\begin{pmatrix} FV_1 & FV_2\\ uS_0 & dS_0\end{pmatrix}
\]
\item Find necessary and sufficient conditions for the model to be
complete.
\item Show that this model is arbitrage--free if and only if
\begin{equation}
\min\left(\frac{u}{FV_1}, \frac{d}{FV_2}\right) \;<\; 1 \;<\;
\max\left(\frac{u}{FV_1}, \frac{d}{FV_2}\right).
\end{equation}
\item Show that, if $FV_1 = FV_2 = e^{r\delta t}$, the condition (3.50) is
equivalent to the no--arbitrage condition $d < e^{r\delta t} < u$ for the
classical one period binomial model.
\end{subparts}

\item In a one period trinomial model, it is assumed that the price at
time $\tau > t_0$ of an asset with price $S_0$ at time $t_0$ will be
either $dS_0$, $mS_0$, or $uS_0$, where $d < m < u$.

Consider a one period market model with two securities, i.e., the asset
described above and cash, and three states, i.e., state $\omega^1$, when
the price of the asset at time $\tau$ is $dS_0$, state $\omega^2$, when the
price of the asset at time $\tau$ is $mS_0$, and state $\omega^3$, when the
price of the asset at time $\tau$ is $uS_0$.
\begin{subparts}
\item Show that the one period trinomial model is incomplete.
\item Show that the one period trinomial model is arbitrage--free if and
only if
\[
d \;<\; e^{r\delta t} \;<\; u,
\]
where $\delta t = \tau - t_0$.
\end{subparts}

\item In three months, the value of an asset with spot price \$40 will be
either \$32, \$38, \$42, or \$44. The value of a three months European
call option with strike \$36 on this asset is \$8 and the value of a three
months European put option with strike \$40 on this asset is \$5. For
simplicity, assume zero risk free rates. Consider the one period market
model with the following four securities and the following four states in
three months:

\noindent Securities:
\begin{itemize}
\item cash;
\item asset;
\item three months call with strike \$36;
\item three months put with strike \$40;
\end{itemize}

\noindent Market states:
\begin{itemize}
\item asset price \$32;
\item asset price \$38;
\item asset price \$42;
\item asset price \$44.
\end{itemize}
\begin{subparts}
\item Find the payoff matrix of this model, and show that the four
securities are non--redundant.
\item Show that the one period market model is complete.
\item How do you replicate a bull spread made of a long position in a
three months call with strike \$34 and a short position in a three months
call with strike \$40 in this one period market?
\item Show that this model is not arbitrage--free, and find an arbitrage
opportunity.
\end{subparts}

\item Two assets have spot prices \$20 and \$30, respectively. Assume
that, in five months, the first asset will be worth either \$18 or \$22,
and the second asset will be worth either \$28 or \$32. Also, assume that
the risk--free rate for five months cash deposits is $0$. Consider the one
period market model with the following three securities and the following
four states in five months:

\noindent Securities:
\begin{itemize}
\item cash;
\item first asset;
\item second asset.
\end{itemize}

\noindent Market states:
\begin{itemize}
\item first asset at \$22 and second asset at \$32;
\item first asset at \$22 and second asset at \$28;
\item first asset at \$18 and second asset at \$32;
\item first asset at \$18 and second asset at \$28.
\end{itemize}
\begin{subparts}
\item Show that this one period market model is arbitrage--free.
\item Show that, in this one period market model, it is not possible to
replicate a derivative security that pays \$1 if the first state occurs
(i.e., if, in five months, the first asset is worth \$22 and the second
asset is worth \$32), and does not pay anything if any other state occurs.
Conclude that the model is not complete.
\end{subparts}
\noindent Hint: For (i), consider the case when all state prices are equal
to $\frac{1}{4}$.

\item This exercise is related to the example from Chapter 1.

Consider two assets with spot prices \$30 and \$50, respectively. Assume
that, in three months, the first asset will be worth either \$34 or \$24,
and the second asset will be worth either \$56, \$51, or \$46. The value
of a three months at--the--money European call option with strike \$30 on
the first asset is \$2.5, the value of a three months at--the--money
European call option with strike \$50 on the second asset is \$2.7, and
the value of a three months European put option with strike \$52 on the
second asset is \$4.1. Assume that the future value in three months of \$1
today is \$1.01. Consider the one period market model with the following
six securities and the following six states in three months:

\noindent Securities:
\begin{itemize}
\item cash;
\item first asset;
\item second asset;
\item three months ATM call on the first asset;
\item three months ATM call on the second asset;
\item three months put with strike \$52 on the second asset;
\end{itemize}

\noindent Market states:
\begin{itemize}
\item first asset at \$34 and second asset at \$56;
\item first asset at \$34 and second asset at \$51;
\item first asset at \$34 and second asset at \$46;
\item first asset at \$24 and second asset at \$56;
\item first asset at \$24 and second asset at \$51;
\item first asset at \$24 and second asset at \$46.
\end{itemize}
\begin{subparts}
\item Show that the payoff matrix for this market model is nonsingular,
and conclude that the market is complete.
\item Compute the state prices for this model, and show that the market
is arbitrage--free.
\item Use risk--neutral pricing to find the value of a three months call
option with strike \$52 on the second asset.
\item Use Put--Call parity to find the value of a three months call
option with strike \$52 on the second asset, and compare it to the value
computed at (iii).
\item What is the value of a bear spread made of a long position in a
three months put option with strike \$35 on the first asset and a short
position in a three months put option with strike \$28 on the first
asset?
\end{subparts}

\item This exercise refers to the S\&P 500 options prices from Table 3.1.

Consider a one period market model with the following nine securities:
\begin{center}
P1200; P1275; P1350; P1375; C1375; C1400; C1450; C1550; C1600.
\end{center}
The nine states of the index price at maturity are as follows: seven
states correspond to the midpoints between the strikes of the options
above, i.e.,
\begin{align*}
\omega^2 &: \{S(\tau) = 1237.50\}; & \omega^6 &: \{S(\tau) = 1425\}; \\
\omega^3 &: \{S(\tau) = 1312.50\}; & \omega^7 &: \{S(\tau) = 1500\}; \\
\omega^4 &: \{S(\tau) = 1362.50\}; & \omega^8 &: \{S(\tau) = 1575\}; \\
\omega^5 &: \{S(\tau) = 1387.50\}; & &
\end{align*}
the first and last state are
\[
\omega^1 \;:\; \{S(\tau) = 950\}; \qquad \omega^9 \;:\; \{S(\tau) = 1675\}.
\]
\begin{subparts}
\item Find the payoff matrix $M_\tau$ of this one period market model, and
show that the model is complete.
\item Find the state prices vector $Q$ and show that the model is
arbitrage--free.
\item Compute the root--mean--squared error (RMSE) of this model. Comment
on the precision of this nine securities model compared to the seven
securities model from section 3.5.
\end{subparts}

\item This problem refers to the S\&P 500 options prices from Table 3.1.

Consider a one period market model with the same seven securities as in
section 3.5, i.e.,
\begin{center}
P1200; P1300; P1400; C1400; C1450; C1550; C1600.
\end{center}
The states of this market are the midpoints of the strike, i.e.,
\begin{align*}
\omega^2 &: \{S(\tau) = 1250\}; & \omega^5 &: \{S(\tau) = 1500\}; \\
\omega^3 &: \{S(\tau) = 1350\}; & \omega^6 &: \{S(\tau) = 1575\}; \\
\omega^4 &: \{S(\tau) = 1425\}; & &
\end{align*}
and the first and last state are
\[
\omega^1 \;:\; \{S(\tau) = 1100\}; \qquad \omega^7 \;:\; \{S(\tau) = 1700\}.
\]
Show that this market model is not arbitrage--free.

\end{exercises}

% PDF page 126 (the last page of this assignment) is blank -- a
% transitional blank verso before Chapter 4, which begins on the next
% recto page. Nothing to transcribe on that page. This is also the end of
% Chapter 3 "The Arrow--Debreu one period market model."
